\chapter*{Предисловие}
\addcontentsline{toc}{chapter}{Предисловие}
\markboth{Предисловие}{}

Аппаратура управления --- это «руль» любой радиоуправляемой модели: квадрокоптера, самолёта,
планера, вертолёта, машины или лодки. Пульты RadioMaster TX12 и Boxer работают под управлением
открытой прошивки \textbf{EdgeTX}, которая превращает их в очень гибкий программируемый компьютер.
С одной стороны, это даёт огромную свободу: практически любую логику управления можно собрать
из стандартных «кубиков». С другой --- у новичка глаза разбегаются от десятка меню, где
всё связано со всем.

Задача этой книжки --- провести вас от распаковки пульта до уверенной самостоятельной настройки
любой модели. Мы не будем пересказывать каждый пункт меню подряд; вместо этого разберём,
\emph{как течёт сигнал} от стика до сервомашинки, и на каждом шаге покажем, какие настройки
на этот сигнал влияют и зачем они нужны.

\section*{Как читать}

\begin{itemize}
  \item Главы 1--4 --- обязательный минимум: устройство пультов, что такое EdgeTX, обновление
        прошивки и системные настройки. Прочитайте их по порядку.
  \item Главы 5--6 --- создание модели и привязка приёмника. После них модель уже можно поднять
        в воздух (или выехать на трассу).
  \item Главы 7--12 --- «сердце» EdgeTX: входы, миксеры, выходы, кривые, полётные режимы,
        глобальные переменные, логические переключатели и специальные функции.
  \item Главы 13--15 --- телеметрия, Lua-скрипты, тренерский режим и симуляторы.
  \item Глава 16 --- готовые рецепты для типовых моделей, глава 17 --- диагностика неполадок.
  \item Приложения --- шпаргалка по кнопкам, структура SD-карты и словарь терминов.
\end{itemize}

\section*{Соглашения}

\begin{itemize}
  \item Физические кнопки пульта обозначены так: \btn{SYS}, \btn{MDL}, \btn{TELE}, \btn{RTN},
        \btn{PAGE>}, \btn{ENTER}.
  \item Пункты меню и параметры набраны \menu{таким шрифтом}. Путь по меню записывается
        через треугольник: \mpath{SYS\sep Hardware\sep Calibration}.
  \item Переключатели и источники сигнала выделены: \sw{SA}, \sw{S1}, \sw{CH5}, \sw{L1}.
  \item EdgeTX можно собрать с русским интерфейсом, но большинство инструкций, видео и форумов
        используют английские названия. Поэтому в книге указано \emph{английское} название
        пункта меню и, при первом упоминании, его перевод.
\end{itemize}

\begin{caution}
Радиоуправляемая модель с работающими пропеллерами опасна. Все настройки, связанные
с двигателем (арм, отсечка газа, failsafe), проверяйте \textbf{со снятыми пропеллерами}.
Не летайте над людьми и соблюдайте правила использования воздушного пространства и
радиочастот, действующие в вашей стране.
\end{caution}

\begin{note}
EdgeTX активно развивается: названия некоторых пунктов и их расположение от версии к версии
немного меняются. Если что-то на вашем экране выглядит иначе --- ищите пункт с похожим
названием на соседних страницах, а подробности --- в официальной документации
\code{edgetx.org} и в заметках к выпуску (release notes) вашей версии.
\end{note}
